\section{Conclusion}
\label{sec:conclusion}

\system\ is a practical multi-provider red-teaming framework with traceable
campaign execution and publication-grade reporting outputs. The refreshed
artifact set in this paper covers a fully reachable curated registry of 11
models, each evaluated on the same 20-probe campaign.

The resulting picture is operationally useful: model families differ in hit
rate, aggregate risk, and execution stability, and those dimensions should be
reported jointly. The framework's value is not only attack discovery but also
transparent, reproducible accounting of the conditions under which findings were
obtained.

Methodologically, continuous red teaming should remain adaptive in both attack
technique and risk scope. Practical programs should include not only obvious
high-severity jailbreak targets but also domains often deprioritized as
``lower-risk,'' since those interactions can still reveal behaviors that matter
for downstream safety evaluations and governance.

Future work should add evaluator-judge reruns over the same curated cohort,
expand curated coverage while preserving reachability guarantees, and execute
engine-native recon/confirm workflows end-to-end. Priority extensions include
multi-judge evaluator ensembles (for example combining frontier API judges with
open-weight judges to stress evaluator sensitivity) and agentic multi-agent or
swarm scenarios in which coordinated tool use, delegation, and memory-sharing
create failure modes that single-agent prompt tests can miss.
